% =============================================================================
% SECTION 1: FOUNDATIONAL AXIOMS AND THE PARTITION POSTULATE
% =============================================================================

\section{Foundational Structure: The Partition Postulate}
\label{sec:axioms}

\subsection{The Problem of Individuation}
\label{subsec:individuation}

Physics presupposes the existence of distinct objects. Newton's laws speak of
particles; thermodynamics speaks of systems and environments; quantum mechanics
speaks of states and observables. In each case, the objects whose behaviour is
to be described are assumed to be already individuated --- already distinct from
one another and from the background against which they move. The question of
\emph{how} individuation arises is not addressed. It is taken as given.

This assumption is not innocent. The structure of every physical law depends on
which items are considered distinct, and the choice of what counts as a system
versus an environment is not determined by the formalism itself. We argue here
that individuation is not a pre-physical given but a physical process with
definite structure, finite duration, and irreducible residue. Formalising
individuation is the first task of this paper.

\subsection{Existence by Negation}
\label{subsec:negation}

\begin{postulate}[Existence by Negation]
\label{post:negation}
An item $\mathcal{I}$ exists as a distinct entity if and only if there is a
well-defined complement $\overline{\mathcal{I}}$ --- the rest of reality --- such
that $\mathcal{I} \cap \overline{\mathcal{I}} = \emptyset$ and
$\mathcal{I} \cup \overline{\mathcal{I}} = \mathcal{U}$, where $\mathcal{U}$
denotes the totality of reality. The existence of $\mathcal{I}$ is constituted
by the distinction $\mathcal{I} \neq \overline{\mathcal{I}}$, not by any
intrinsic property of $\mathcal{I}$ alone.
\end{postulate}

This postulate makes explicit what is implicit in all of physics: to speak of an
item is to speak of something that is \emph{not} everything else. The item has no
independent existence prior to this distinction. Its properties --- mass, charge,
position, momentum --- are relational consequences of the distinction, not
pre-existing attributes that the distinction then separates from the background.

\begin{remark}
Postulate~\ref{post:negation} does not presuppose a spatial background. The
complement $\overline{\mathcal{I}}$ is not ``the rest of space'' but ``the rest
of reality''. Space, as we show in Section~\ref{sec:emergence}, is itself an
emergent consequence of partition structure, not the container within which
partitions occur.
\end{remark}

\subsection{The Partition Process}
\label{subsec:partition_process}

The distinction $\mathcal{I} \neq \overline{\mathcal{I}}$ does not arise
instantaneously. If it did, there would be no difference between the state of
the universe before and after the distinction came into being --- which would
mean the distinction made no difference, which would mean $\mathcal{I}$ had not
actually become distinct. We therefore require:

\begin{postulate}[Partition Duration]
\label{post:duration}
The establishment of the distinction $\mathcal{I} \neq \overline{\mathcal{I}}$
requires a finite, non-zero duration $\tau_p > 0$, called the
\emph{partition interval}. During $\tau_p$, the universe $\mathcal{U}$ must
progress --- that is, the state of $\overline{\mathcal{I}}$ must change --- for
the partition to be well-defined at its completion. If $\overline{\mathcal{I}}$
did not change during $\tau_p$, then the partition would have occurred in a
frozen universe, and the item $\mathcal{I}$ would be indistinguishable from the
state $\mathcal{U}$ had before the partition began.
\end{postulate}

\begin{definition}[Partition State]
\label{def:partition_state}
A \emph{partition state} is a quadruple $(n, \ell, m, s)$ where:
\begin{itemize}
    \item $n \in \mathbb{Z}^+$ is the \emph{partition depth} --- the number of
          nested boundary layers separating $\mathcal{I}$ from $\mathcal{U}$
    \item $\ell \in \{0, 1, \ldots, n-1\}$ is the \emph{structural complexity}
          --- the number of internal sub-partitions of $\mathcal{I}$
    \item $m \in \{-\ell, \ldots, +\ell\}$ is the \emph{orientation index}
          --- the asymmetry of the boundary relative to $\overline{\mathcal{I}}$
    \item $s \in \{-\frac{1}{2}, +\frac{1}{2}\}$ is the \emph{residue parity}
          --- the sign of the irreducible residue generated by the partition
\end{itemize}
The set of all partition states is denoted $\mathcal{P}$.
\end{definition}

\begin{theorem}[Partition Capacity]
\label{thm:capacity}
The number of distinct partition states at depth $n$ is $C(n) = 2n^2$. The
cumulative number of partition states accessible at depth $\leq n$ is:
\begin{equation}
    N_{\mathrm{state}}(n) = \sum_{k=1}^{n} C(k) = \frac{n(n+1)(2n+1)}{3}
    \label{eq:cumulative_states}
\end{equation}
\end{theorem}

\begin{proof}
At depth $n$, the structural complexity $\ell$ ranges over $n$ values
$\{0, 1, \ldots, n-1\}$. For each $\ell$, the orientation index $m$ takes
$2\ell + 1$ values. The residue parity $s$ takes 2 values independently of
$\ell$ and $m$. Therefore:
\begin{equation}
    C(n) = 2 \sum_{\ell=0}^{n-1}(2\ell + 1) = 2n^2
\end{equation}
The cumulative count follows by summing $C(k) = 2k^2$ from $k=1$ to $n$,
using $\sum_{k=1}^n k^2 = n(n+1)(2n+1)/6$.
\end{proof}

\subsection{The Bijection and the Partition Count}
\label{subsec:bijection}

\begin{theorem}[Partition Bijection]
\label{thm:bijection}
There exists a canonical bijection $\Phi: \mathbb{Z}^+ \to \mathcal{P}$ that
maps each positive integer $M$ to a unique partition state $(n,\ell,m,s)$, where
$n$ is determined by $N_{\mathrm{state}}(n-1) < M \leq N_{\mathrm{state}}(n)$.
\end{theorem}

\begin{proof}
The map $\Phi$ is constructed by ordering partition states lexicographically:
first by $n$ (ascending), then by $\ell$ (ascending), then by $m$ (ascending),
then by $s$ (ascending). Since $N_{\mathrm{state}}(n) = n(n+1)(2n+1)/3$ is
strictly increasing in $n$ and $\mathcal{P} = \bigcup_{n=1}^\infty
\mathcal{P}_n$ where $|\mathcal{P}_n| = C(n) = 2n^2$, the ordering provides a
bijection between $\mathbb{Z}^+$ and $\mathcal{P}$.
\end{proof}

The integer $M = \Phi^{-1}(n,\ell,m,s)$ is the \emph{partition count} --- the
cumulative number of partition events that have occurred in the history of the
system. It is the fundamental observable of the framework.

\subsection{The Irreducible Residue}
\label{subsec:residue}

\begin{postulate}[Residue Irreducibility]
\label{post:floor}
Every partition generates a residue $\beta > 0$ --- a non-zero remainder of the
partition process that cannot be eliminated. The residue satisfies:
\begin{equation}
    \beta \geq \beta_{\min} > 0
    \label{eq:floor}
\end{equation}
where $\beta_{\min}$ is the \emph{floor} of the residue, dependent only on the
partition depth $n$ and the partition interval $\tau_p$.
\end{postulate}

The necessity of $\beta > 0$ follows from Postulate~\ref{post:duration}. During
the partition interval $\tau_p$, the complement $\overline{\mathcal{I}}$ has
progressed. When the partition completes, the boundary between $\mathcal{I}$ and
$\overline{\mathcal{I}}$ therefore has a \emph{thickness} --- it is not a
perfect cut but a layer of states that are neither fully $\mathcal{I}$ nor fully
$\overline{\mathcal{I}}$. This thickness is $\beta$.

\begin{corollary}[No Perfect Partition]
\label{cor:noperfect}
No partition can achieve $\beta = 0$. A partition with $\beta = 0$ would require
either zero partition duration ($\tau_p = 0$, violating
Postulate~\ref{post:duration}) or a frozen universe during the partition (also
violating Postulate~\ref{post:duration}).
\end{corollary}

\subsection{Complete Knowledge and the Boundary Thickness}
\label{subsec:knowledge}

The impossibility of $\beta = 0$ has an epistemological consequence that is
also a physical consequence. To have complete knowledge of $\mathcal{I}$ would
require knowing everything about $\mathcal{I}$ and nothing about
$\overline{\mathcal{I}}$, or vice versa. But the boundary has thickness $\beta$
--- a layer of states that belong to neither. Any agent with complete knowledge
of $\mathcal{I}$ would have to have complete knowledge of the boundary, which
means partial knowledge of $\overline{\mathcal{I}}$. And any agent with complete
knowledge of $\overline{\mathcal{I}}$ would have to have complete knowledge of
the boundary, which means partial knowledge of $\mathcal{I}$.

The boundary thickness $\beta$ is therefore both:
\begin{enumerate}
    \item A physical quantity --- the irreducible layer of states at the
          partition boundary
    \item An epistemological bound --- the irreducible uncertainty in the
          distinction between $\mathcal{I}$ and $\overline{\mathcal{I}}$
\end{enumerate}

These two are not two descriptions of the same thing. They are the same thing.

\begin{theorem}[Floor Theorem]
\label{thm:floor}
For a system undergoing repeated partitioning, the partition count $M$ satisfies:
\begin{equation}
    M(t_2) > M(t_1) \quad \text{for all } t_2 > t_1
    \label{eq:monotone}
\end{equation}
That is, the partition count is strictly monotone increasing. No physical process
can decrement $M$.
\end{theorem}

\begin{proof}
Each partition event increments $M$ by exactly 1 (the bijection
$\Phi: \mathbb{Z}^+ \to \mathcal{P}$ maps each new partition state to the next
integer). A partition event requires $\tau_p > 0$ (Postulate~\ref{post:duration})
and generates residue $\beta > 0$ (Postulate~\ref{post:floor}). The residue
$\beta$ is deposited into the contact chain of $\mathcal{I}$ with
$\overline{\mathcal{I}}$. For $M$ to decrease, a partition event would have to
``un-occur'' --- the residue $\beta$ would have to be withdrawn from the contact
chain and the distinction $\mathcal{I} \neq \overline{\mathcal{I}}$ would have
to be un-established. But un-establishing the distinction would require
$\mathcal{I}$ to re-merge with $\overline{\mathcal{I}}$, which requires all
contact relationships established by the partition to simultaneously dissolve.
Since each contact relationship has itself generated a residue (by
Postulate~\ref{post:floor}), dissolving it requires the same argument
recursively, expanding without bound. Therefore $M$ cannot decrease.
\end{proof}
